\documentclass[11pt]{article}
\usepackage[margin=0.84in]{geometry}
\usepackage{amsmath,amssymb,mathtools,booktabs,microtype,xcolor,fancyhdr,hyperref}
\definecolor{navy}{HTML}{17365D}\definecolor{teal}{HTML}{147D78}
\hypersetup{colorlinks=true,linkcolor=navy,urlcolor=teal}
\pagestyle{fancy}\setlength{\headheight}{14pt}\fancyhf{}
\lhead{Ising/Jones $B_3$ image}\rhead{Proof and verification}\cfoot{\thepage}
\newcommand{\Sgm}{\mathcal S}
\title{\color{navy}\bfseries The Finite Ising/Jones $B_3$ Image\\
\large Lift-dependent and projectively intrinsic sigma--MGFs}
\author{Computational verification note}\date{17 July 2026}
\begin{document}\maketitle

\begin{abstract}
We test the finite-image formulas on the two-dimensional Ising/Jones braid
representation.  The displayed braid generators close to a 192-element
unitary lift with scalar center of order eight.  We verify its one-sided
sigma--MGF and the full-twist selection rule.  We then pass to conjugation on
$\operatorname{End}(\mathbb C^2)$, obtaining the canonical 24-element
projective image in dimension four, and verify the projective formula there.
\end{abstract}

\section{The finite braid image}
Put
\[
 H=\frac1{\sqrt2}\begin{pmatrix}1&1\\1&-1\end{pmatrix},\qquad
 S=\begin{pmatrix}1&0\\0&i\end{pmatrix},\qquad q=e^{-\pi i/8},
\]
and define
\[
 \rho(\sigma_1)=qS,\qquad \rho(\sigma_2)=qHSH.
\]
Direct multiplication gives the braid relation.  Breadth-first matrix closure
gives $|G|=192$; quotienting its eight scalar matrices gives a projective image
of order $24$.  The $B_3$ full twist is
\[
 \rho((\sigma_1\sigma_2)^3)=e^{-\pi i/4}I.
\]
It therefore acts on
$W_\tau=\bigotimes_a\operatorname{Sym}^{\tau_a}\mathbb C^2$ by
$e^{-\pi i|\tau|/4}$, proving the necessary selection rule
\[
 \boxed{W_\tau^G=0\quad\text{unless}\quad8\mid|\tau|.}
\]
The rule is necessary, not sufficient, as the degree-eight computations show.

\section{Finite-image Molien proof}
For every finite unitary image $G\leq U(V)$,
\[
 \Sgm_{G,V}(\mathbf t)=\frac1{|G|}\sum_{g\in G}
 \prod_{a\geq1}\det(I-t_ag)^{-1}.
\]
Since the coefficient of $t^r$ is the character of
$\operatorname{Sym}^rV$, Reynolds averaging proves
\[
 [m_\tau]\Sgm_{G,V}=\dim(W_\tau^G).
\]
Taking a logarithm gives
\[
 \det(I-tg)^{-1}=\exp\left(\sum_{k\geq1}
 \frac{t^k}{k}\operatorname{Tr}(g^k)\right),
\]
and partitioning the sum into conjugacy classes proves the character/power-map
and centralizer-weighted class formulas.  No information about the infinite
kernel of $B_3\to G$ is needed after the finite image is fixed.

\section{Projective correction}
Changing a lift $g$ to $zg$ changes its eigenvalues, so the one-sided series
is not a projective invariant.  Conjugation on $\operatorname{End}(V)$ is
unchanged:
\[
 \operatorname{Ad}_{zg}(A)=(zg)A(zg)^{-1}=gAg^{-1}.
\]
Thus the projectively intrinsic series is the ordinary finite Molien series
of the adjoint representation,
\[
 \Sgm^{\mathrm{proj}}(\mathbf t)=\frac1{|\overline G|}
 \sum_{\bar g\in\overline G}\prod_a
 \det(I-t_a\operatorname{Ad}_{\bar g})^{-1}.
\]
If $g$ is any lift, then
\[
 \operatorname{Tr}(\operatorname{Ad}_{g}^{,k})
 =\operatorname{Tr}(g^k)\overline{\operatorname{Tr}(g^k)}
 =|\chi(g^k)|^2,
\]
which proves the proposed projective power-trace formula and its independence
of lift.

\section{Independent matrix checks}
For each $\tau$, the program explicitly restricts tensor powers to normalized
symmetric occupation bases and ranks the averaged Reynolds matrix.  Separately,
it computes $h_r(g)$ from traces of powers using Newton's recurrence and
averages $\prod_a h_{\tau_a}(g)$ first over elements and then over matrix
conjugacy classes.  The lift has 32 classes; the adjoint image has five.

\begin{center}
\begin{tabular}{@{}llrrrr@{}}\toprule
representation&$\tau$&$\dim W_\tau$&rank&elements&classes\\\midrule
lift on $\mathbb C^2$&$(1)$&2&0&0&0\\
&$(4)$&5&0&0&0\\
&$(8)$&9&0&0&0\\
&$(4,4)$&25&1&1&1\\\addlinespace
adjoint on $\operatorname{End}(\mathbb C^2)$&$(1)$&4&1&1&1\\
&$(2)$&10&2&2&2\\
&$(1,1)$&16&2&2&2\\
&$(3)$&20&2&2&2\\\bottomrule
\end{tabular}
\end{center}
All braid, unitarity, closure, integrality, and equality assertions pass.  The
largest projector error is below $4\times10^{-15}$.

\section{Reproduction and interpretation}
Run
\begin{verbatim}
.venv/bin/python lambda_distributions/proofs/finite_braid_images/
  ising_clifford/verification.py
.venv/bin/marimo edit marimo_notebooks/braid_ising_clifford.py
\end{verbatim}
after joining wrapped lines.  The computation verifies the stated formulas
for the displayed $B_3$ Ising representation.  It does not claim that every
root-of-unity Jones representation has finite image; finiteness remains an
essential hypothesis.
\end{document}

